\section{The Capstone: The Model Run End to End}
\label{sec:capstone}

The committed model is not an assembly of separate experiments. It is one object, and the testbed
runs it end to end (P34, P36, P52 to P56), reading the key emergent structures off a single
instantiation.

\paragraph{One mesh, many structures (P34).}
P34 is the capstone run. A single committed mesh is built and updated, and the emergent structures
are read off the same instance: the mean degree and exponential reach (the hyperbolic geometry), the
Lorentz anisotropy (consistent with zero), the bounded-below Hamiltonian (the energy), the
integration (the unity), and the tone histogram. They come from one mesh, not from separate
hand-built models, which is the real test of unification.

\paragraph{The model constructor (P36).}
P36 is the small constructor, \texttt{code/model/vibe.ts}, that writes the committed model at a
glance and exposes the substrate choices: the hyperbolic mesh, the heptagrid and the dodecagrid as
the ideal crystals it approximates, and the lattice and sprinkle as comparison meshes. Its
\texttt{read()} returns the emergent structures above, including the integration and a count of the
higher vibes coarse-grained off the settled mesh.

\paragraph{Hardening (P52 to P54).}
P52 to P54 are the hardening pass. P52 checks the continuum limit, with the emergent dimension
agreeing across sizes. P53 finds a coarse-graining fixed point, the macro-rule reproducing the
micro-rule under renormalization. P54 measures large-$N$ scaling so the heavier experiments are
trustworthy at size.

\paragraph{The unified operator and the eternal ladder (P55, P56).}
P55 brings all the bosonic sectors out of a single operator on a single mesh, the unification at
the level of the substrate. P56 sets the integer ladder of Section~\ref{sec:base} growing eternally
with the model on it, joining the static base to the deterministic eternal growth of
Section~\ref{sec:geometry} (P83). The crystal, the rule, and the growth are one system.
